\input{/Users/etiennecollin/Programming/LaTex/Paths}
\path{macOS}

\documentclass[12pt]{article}
\usepackage[english]{babel}																												% Selects language of document
\usepackage{./_assets/Packages}																											% Imports custom packages selection
\usepackage[style=numeric,autocite=superscript,sorting=none,backref=false,backend=biber,date=iso,datezeros=true,seconds=true]{biblatex}	% Bibliography options. autocite=<inline, footnote, superscript, plain>
% \usepackage{indentfirst}																												% Used to indent first paragraph of section

% DOCUMENT USER SETTINGS ==============================================================================================
\newcommand{\docAuthorName}{Etienne Collin, Marie Ouellet \& Rania Yahyaoui}
\newcommand{\docAuthorStudentNumber}{}
\newcommand{\docAuthorTitlePage}{\docAuthorName}
\newcommand{\docClass}{Probability and Statistics}
\newcommand{\docClassInstructor}{Professor Vincent Carrier}
\newcommand{\docClassNumber}{201-BNM-LW}
\newcommand{\docClassTime}{Section 20108}
\newcommand{\docClassSemester}{Winter 2022}
\newcommand{\docDueDate}{May 25, 2022}
\newcommand{\docDueTime}{23:59}
\newcommand{\docSubTitle}{Properties, Proofs and Applications}
\newcommand{\docTitle}{The Gamma Distribution}
\usepackage[svgnames, hyperref]{xcolor}	% See p.42 https://mirror.its.dal.ca/ctan/macros/latex/contrib/xcolor/xcolor.pdf

% Page settings
\topmargin=-0.45in
\evensidemargin=0in
\oddsidemargin=0in
\textwidth=6.5in
\textheight=9.0in
\headsep=0.25in
\linespread{1.1}
\pagestyle{fancy}
\fancyhf[HL]{\firstxmark}						% Content on left of head
\fancyhf[HC]{\docTitle}						% Content on center of head
\fancyhf[HR]{Collin - \docAuthorStudentNumber}			% Content on right of head
\fancyhf[FL]{\docClassInstructor}					% Content on left of foot
\fancyhf[FC]{\thepage}						% Content on center of foot
\fancyhf[FR]{\lastxmark}						% Content on right of foot
\renewcommand\headrulewidth{0.4pt}				% Width of line separation head from body
\renewcommand\footrulewidth{0.4pt}				% Width of line separation foot from body
% \setlength\parindent{5mm}					% Custom Paragraph indentation
\setlength\parindent{0.5in}					% Custom Paragraph indentation											% Imports custom page settings
% DEFINITIONS =========================================================================================================
% Paths
\addbibresource{\basepath/bibliography.bib}	% Imports bibliography file
\graphicspath{{images/}}

% lstlistings setup ---------------------------------------------------------------------------------------------------
\definecolor{commentsColor}{gray}{0.6}
\definecolor{keywordsColor}{RGB}{249,38,114}
\definecolor{stringColor}{RGB}{230, 219, 116}
\definecolor{mygreen}{rgb}{0,0.6,0}
\definecolor{mygray}{rgb}{0.5,0.5,0.5}
\definecolor{mymauve}{rgb}{0.58,0,0.82}
\lstset{
  backgroundcolor=\color{white},   % choose the background color; you must add \usepackage{color} or \usepackage{xcolor}; should come as last argument
  basicstyle=\footnotesize,        % the size of the fonts that are used for the code
  breakatwhitespace=false,         % sets if automatic breaks should only happen at whitespace
  breaklines=true,                 % sets automatic line breaking
  captionpos=b,                    % sets the caption-position to bottom
  commentstyle=\color{mygreen},    % comment style
  deletekeywords={},            % if you want to delete keywords from the given language
  escapeinside={\%*}{*)},          % if you want to add LaTeX within your code
  extendedchars=true,              % lets you use non-ASCII characters; for 8-bits encodings only, does not work with UTF-8
  firstnumber=0,                % start line enumeration with line 0
  frame=single,	                   % adds a frame around the code
  keepspaces=true,                 % keeps spaces in text, useful for keeping indentation of code (possibly needs columns=flexible)
  keywordstyle=\color{blue},       % keyword style
  language=Python,                 % the language of the code
  morekeywords={*},            % if you want to add more keywords to the set
  numbers=left,                    % where to put the line-numbers; possible values are (none, left, right)
  numbersep=5pt,                   % how far the line-numbers are from the code
  numberstyle=\tiny\color{mygray}, % the style that is used for the line-numbers
  rulecolor=\color{black},         % if not set, the frame-color may be changed on line-breaks within not-black text (e.g. comments (green here))
  showspaces=false,                % show spaces everywhere adding particular underscores; it overrides 'showstringspaces'
  showstringspaces=false,          % underline spaces within strings only
  showtabs=false,                  % show tabs within strings adding particular underscores
  stepnumber=2,                    % the step between two line-numbers. If it's 1, each line will be numbered
  stringstyle=\color{mymauve},     % string literal style
  tabsize=4,	                   % sets default tabsize to 2 spaces
  title=\lstname                   % show the filename of files included with \lstinputlisting; also try caption instead of title
}

% Sections alignment --------------------------------------------------------------------------------------------------
% Gets label of current section
\makeatletter\newcommand{\currentlabel}{\@currentlabelname}\makeatother
% Sets alignment of sections
\newcommand{\formatsection}[1]{
    \ifthenelse{\equal{#1}{center}}{
		%Makes sections centered
		\titleformat{\section}{\normalfont\Large\bfseries\centering}{\thesection}{1em}{}
    }{
        \ifthenelse{\equal{#1}{right}}{
			%Makes sections right-aligned
			\titleformat{\section}{\normalfont\Large\bfseries\raggedleft}{\thesection}{1em}{}
        }{
			\ifthenelse{\equal{#1}{left}}{
				%Makes sections left-aligned
				\titleformat{\section}{\normalfont\Large\bfseries\raggedright}{\thesection}{1em}{}
        	}{
        	}
        }
    }
}

% tabularx
\newcolumntype{Y}{>{\centering\arraybackslash}X}

% Caption/Ref settings ----------------------------------------------------------------------------------------------------
\captionsetup{
    format = hang,
    labelfont = sc,
	skip = 6pt,
	hypcapspace = 96pt
}

% TOC -----------------------------------------------------------------------------------------------------------------
% Sets TOC depth and change changes its display name
\setcounter{secnumdepth}{3} % Change to 0 for unnumbered sections
\addto\captionsenglish{
  \renewcommand{\contentsname}{Table of Contents}
}

% SIunitx -------------------------------------------------------------------------------------------------------------
% Sets \per to display as frac
\sisetup{per-mode=symbol-or-fraction}
\DeclareSIUnit\Molar{\mole\per\liter}
\DeclareSIUnit\molar{\textsc{M}}
\DeclareSIUnit\clight{\text{\ensuremath{c}}_{0}}

% BibLaTex ------------------------------------------------------------------------------------------------------------
% BibLaTex definitions
\let\oldcite\cite
\renewcommand{\cite}{\autocite}
\let\oldsupercite\supercite
\renewcommand{\supercite}[1]{\textsuperscript{[}\oldsupercite{#1}\textsuperscript{]}}
% Makes number created by \supercite to be between brackets.
\DeclareCiteCommand{\supercite}[\mkbibsuperscript]{%
	\iffieldundef{prenote}{}{\BibliographyWarning{Ignoring prenote argument}}%
	\iffieldundef{postnote}{}{\BibliographyWarning{Ignoring postnote argument}}%
}
{\usebibmacro{citeindex}\,\bibopenbracket\usebibmacro{cite}\bibclosebracket}
{\supercitedelim}{}

% BASIC DOCUMENT SETTINGS AND TITLE PAGE ==============================================================================
% PDF export settings
\hypersetup{
	pdftex,
	pdfauthor = {Etienne Collin},
	bookmarksnumbered = true,
	bookmarksopen = true,
	bookmarksopenlevel = 1,
	pdfstartview = Fit,
	pdfpagemode = UseOutlines,
	pdfpagelayout = TwoPageRight,
	colorlinks = true,		% Colors links instead of ugly boxes
	urlcolor = DarkBlue,	% Color for external hyperlinks
	linkcolor = black,		% Color of internal links
	citecolor = black,		% Color of citations
	filecolor = black,		% Color of file links
}

% Custom font

\newcommand{\customfont}[1]{
	\usepackage{mathspec}	% Enable the power of CUSTOM FONTS with XeLaTeX
	\setallmainfonts(Digits,Latin,Greek)[Numbers={Lining,Proportional}]{#1}
	\setallmonofonts[Numbers={Lining,Proportional}]{#1}
	\setallsansfonts[Numbers={Lining,Proportional}]{#1}
	\setmathsfont(Digits,Latin,Greek)[Numbers={Lining,Proportional}]{#1}
	\setmathsf[Numbers={Lining,Proportional}]{#1}
	\setmathbb[Numbers={Lining,Proportional}]{#1}
	\setmathcal[Numbers={Lining,Proportional}]{#1}
}

% ENVIRONMENTS ========================================================================================================
% hmwkProblem ---------------------------------------------------------------------------------------------------------
\newcounter{partCounter}
\renewcommand{\part}[1]{\noindent\textbf{\large Part (\alph{partCounter})}\stepcounter{partCounter}\\}
\newcounter{hmwkProblemCounter}
\setcounter{hmwkProblemCounter}{1}

% Continued Header Changer
\newcommand{\enterProblemHeader}[1]{
	\nobreak\extramarks{}{}\nobreak{}
	\nobreak\extramarks{Problem \arabic{#1} (continued)}{Problem \arabic{#1} continued on next page\ldots}\nobreak{}
}
\newcommand{\exitProblemHeader}[1]{
	\nobreak\extramarks{Problem \arabic{#1} (continued)}{Problem \arabic{#1} continued on next page\ldots}\nobreak{}
	\stepcounter{#1}
	\nobreak\extramarks{}{}\nobreak{}
}

% Environment. Optional argument will adjust the problem counter.
\newenvironment{hmwkProblem}[1][-1]{
	\ifnum#1>0
		\setcounter{hmwkProblemCounter}{#1}
	\fi
	\pagebreak
	\section{Problem \arabic{hmwkProblemCounter}}
	\stepcounter{section}
	\setcounter{partCounter}{1}
	\enterProblemHeader{hmwkProblemCounter}
}{
	\exitProblemHeader{hmwkProblemCounter}
}

% sect ----------------------------------------------------------------------------------------------------------------
% Will write "continued" when content spans >1 page. Argument adjusts the name of the section.
\newenvironment{sect}[1]{
	\pagebreak
	\section{#1}
	\stepcounter{section}
	\setcounter{partCounter}{1}
	\nobreak\extramarks{}{}\nobreak{}
	\nobreak\extramarks{\currentlabel\ (continued)}{\currentlabel\ continued on next page\ldots}\nobreak{}
}{
	\nobreak\extramarks{\currentlabel\ (continued)}{\currentlabel\ continued on next page\ldots}\nobreak{}
	\nobreak\extramarks{}{}\nobreak{}
}

% SETUP/ALIASES =======================================================================================================
\newcommand{\imagecenter}[2][1]{
	\begin{center}
		\includegraphics[width=\textwidth*{#1}]{#2}
	\end{center}
}
\newcommand{\imageleft}[2][1]{
	\begin{flushleft}
		\includegraphics[width=\textwidth*{#1}]{#2}
	\end{flushleft}
}
\newcommand{\imageright}[2][1]{
	\begin{flushright}
		\includegraphics[width=\textwidth*{#1}]{#2}
	\end{flushright}
}
\newcommand{\figurecenter}[4]{
	\begin{figure}[H]
		\centering
		\begin{minipage}[H]{#1\textwidth}
			\centering
			\includegraphics[width=\textwidth]{#2}
			\caption{#3}
			\label{#4}
		\end{minipage}
	\end{figure}
}
\newcommand{\figuredual}[6]{
	\begin{figure}[H]
		\begin{minipage}[H]{0.48\textwidth}
			\centering
			\includegraphics[width=\textwidth]{#1}
			\raggedleft
			\caption{#2}
			\label{#3}
		\end{minipage}
		\hfill
		\begin{minipage}[H]{0.48\textwidth}
			\centering
			\includegraphics[width=\textwidth]{#4}
			\caption{#5}
			\label{#6}
		\end{minipage}
	\end{figure}
}
\newcommand{\figuredualsingle}[6]{
	\begin{figure}[H]
		\begin{minipage}[H]{0.48\textwidth}
			\centering
			\includegraphics[width=\textwidth]{#1}
			\subcaption{#2}\label{#6-a}
		\end{minipage}
		\hfill
		\begin{minipage}[H]{0.48\textwidth}
			\centering
			\includegraphics[width=\textwidth]{#3}
			\subcaption{#4}\label{#6-b}
		\end{minipage}
		\caption{#5}\label{#6}
	\end{figure}
}

% Theorems
\theoremstyle{plain}
\newtheorem{definition}{Definition}[section]
\newtheorem{theorem}{Theorem}[section]
\newtheorem{example}{Example}[theorem]
\newtheorem{corollary}{Corollary}[theorem]
\newtheorem{lemma}{Lemma}[theorem]
\newtheorem{proposition}{Proposition}[theorem]
% Unnumbered version
\newtheorem*{definition*}{Definition}
\newtheorem*{theorem*}{Theorem}
\newtheorem*{example*}{Example}
\newtheorem*{corollary*}{Corollary}
\newtheorem*{lemma*}{Lemma}
\newtheorem*{proposition*}{Proposition}

\theoremstyle{definition}
\newtheorem{exercise}{Exercise}[section]
% Unnumbered version
\newtheorem*{exercise*}{Exercise}


\theoremstyle{remark}
\newtheorem{remark}{Remark}
\newtheorem{case}{Case}
\newtheorem{claim}{Claim}[section]
% Unnumbered version
\newtheorem*{remark*}{Remark}
\newtheorem*{case*}{Case}
\newtheorem*{claim*}{Claim}

\renewcommand\qedsymbol{QED}

% Useful shortcuts
\newcommand{\alg}[1]{\textsc{\bfseries\footnotesize #1}}
\newcommand{\deriv}[2][]{\mathop{}\!\frac{\mathrm{d}^{#1}}{\mathrm{d}#2^{#1}}}	% For derivatives
\newcommand{\pderiv}[3][]{\mathop{}\!\dfrac{\partial^{#1}#2}{\partial#3^{#1}}}					% For partial derivatives
\newcommand{\pdiff}[2][]{\mathop{}\!\dfrac{\partial^{#1}}{\partial#2^{#1}}}					% For partial derivatives
\newcommand{\diff}[2][]{\mathop{}\!\mathrm{d^{#1}}#2}							% Integral dx
\newcommand*{\widefbox}[1]{\fbox{\hspace{2em}#1\hspace{2em}}}						% For boxed align environment

% Enables new list environment that supports nesting
\newlist{legal}{enumerate}{10}
\setlist[legal]{label*=.\arabic*}
\setlist[legal,1]{label=\arabic*}

% Temporary text markers
\newcommand{\solution}{\noindent\textbf{\large Solution}}				% Alias for the Solution section header
\newcommand{\missing}{\textcolor{red}{\textbf{VALUE MISSING}}}	% Creates a big red "Missing Value" string.
\newcommand{\inprogress}{\textcolor{red}{\textbf{IN PROGRESS}}}	% Creates a big red "In Progress" string.
\newcommand{\incomplete}{\textcolor{red}{\textbf{INCOMPLETE}}}	% Creates a big red "Incompletes" string.
\newcommand{\bigred}[1]{\textcolor{red}{\textbf{#1}}}	% Creates a big red string.

% END OF SETUP DOCUMENT ===============================================================================================											% Imports custom environments and definitions
\fancyhf[HR]{\docClassTime}													% Removes student number from right header

% SOURCE ==============================================================================================================
\begin{document}
\singlespacing
% Title page
\begin{titlepage}
	\begin{center}
		\vspace*{0.5in}
		\LARGE{\textsc{\docTitle}}\\
		\vspace*{0.1in}
		\Large{\textsc{\docSubTitle}}\\
		\vfill
		%
		\Large{\textsc{\docAuthorTitlePage}}\\
		\vspace*{0.1in}
		\Large{\textsc{\docClass\ -\ \docClassNumber}}\\
		\vfill
		%
		\Large{\docClassTime}\\
		\vspace*{0.1in}
		\large{\textsc{\docClassInstructor}}\\
		\vfill
		%
		\large{\textsc{St. Lawrence Cégep Champlain}}\\
		\vspace*{0.1in}
		\small{Due on \docDueDate\ at \docDueTime}\\
		\vfill
		%
		\includegraphics[width=0.6\textwidth]{../images/Title_page_logo.png}
		\vfill
	\end{center}
\end{titlepage}
\pagebreak
\phantomsection\tableofcontents\pagebreak
\doublespacing
% START OF DOCUMENT ===================================================================================================
\section{General Applications and Family}

The gamma distribution is used in many fields, since it is directly related to the erlang, normal and exponential
distributions whose contributions extend to several disciplines. Here, we will present some of those applications.

First, the gamma distribution is of great use in the field of insurance services, given its direct relation to the
exponential distribution. For example, an analyst could use this distribution to specify the amount of time a product
lasts if one uses it at a constant average rate, thus modeling how reliable it is (and how much insurance should be
charged for it). Then, the size of loan defaults and the cost of insurance claims are also often modeled according to a
gamma distribution.

Then, the gamma distribution can also be used to model the amount of rainfall accumulated in a given reservoir. Indeed,
this distribution fits positive data, represents rainfall distribution well and its two parameters -shape and scale-
give it sufficient flexibility to fit various climates.

Service time can also be modeled using the gamma distribution. For example, if one is waiting in line for a meal, the
waiting time until one receives the long-awaited food can be modelled using an exponential distribution. Using the same
principle, the erlang distribution can allow one to determine the total length of a process, that is, the duration of a
sequence of independent events. For instance, if a large number of people are waiting in line to be served, the
distribution of each of their individual waiting times (the sum of several independent exponentially distributed
variables) will correspond to the time it takes for the employee to serve everyone in it. Therefore, the gamma
distribution lies at the heart of what is called queuing theory: the mathematical study of the congestion of waiting
lines. Since waiting lines are found in countless places such as banks, restaurants and hospitals, as well as on web
servers or multi-step manufacturing and distribution processes, the gamma distribution provides very useful applications
in everyday life.

One of the most famous of those applications concerns phone queuing, on which A. K. Erlang famously worked. The Erlang
distribution has indeed been developed in the goal to model the time in between incoming calls at a call center, along
with the expected number of calls, thus allowing call centers to know what their staffing capacity should be depending
on the time of day.

Beyond waiting lines, the erlang distribution is often used by retailers to model the frequency of interpurchase times
by consumers, which gives them an idea of how often a given consumer is expected to purchase a product from them and
helps them control inventory and staffing.

Finally, in the field of oncology, the age distribution of cancer incidence also follows a gamma distribution. Although
the factors underlying cancer development are not yet fully understood, it has been hypothesized that cancers arise
after several successive “driver events”, that is, after some number of mutations occurs in a cell. Analyses of cancer
statistics suggest that the incidence of the most prevalent cancer types with respect to the patients' age closely
follows the gamma probability distribution or, more specifically, the erlang distribution. This may be due to the fact
that, more broadly, the erlang distribution can be used to model cell cycle time distribution.

\begin{equation}
	Y = A\cdot x^{k-1}\frac{e^{\frac{-x}{b}}}{b^k}\cdot\Gamma(k)
\end{equation}

In this case, the shape parameter \alpha predicts the number of carcinogenic driver events, whereas the shape parameter \lambda 
predicts the average time between those events for each cancer type. Using an additional amplitude parameter A, the
maximal population susceptibility to a given type of cancer can even be predicted. Given that experimental research on
cancer development is crucial for the lives of many people, numerical references such as that provided by the gamma
distribution are of paramount importance in our society. The gamma distribution can save lives, if it is used wisely!

\figurecenter{1}{Cancer_probability}{"The Erlang distribution approximates cancer incidence by age for 20 most prevalent cancer types. Dots indicate actual data for 5-year age intervals, curves indicate the PDF of the Erlang distribution fitted to the data [...]. The middle age of each age group is plotted. Cancer types are arranged in the order of decreasing incidence"\cite{belikovNumberKeyCarcinogenic2017} .}{fig:cancerProbability}
\pagebreak

\section{Introduction}
The gamma distribution is part of the two-parameters family of continuous probability distributions. Indeed, it may be
parameterized with two different parameterizations\cite{wikipediaGammaDistribution2022} :

\noindent Parameterization 1:
\vspace*{-24pt}
\begin{align}\label{eq:gamma:parameterization1}
	\text{Shape: } \alpha>0		&&	\text{Rate: }\lambda>0
\end{align}
Parameterization 2:
\vspace*{-24pt}
\begin{align}\label{eq:gamma:parameterization2}
	\text{Shape: } k>0			&&	\text{Scale: }\theta>0
\end{align}

In the second parameterization, the scale parameter \theta corresponds to the inverse of the rate parameter \lambda, allowing flexibility in the modelization of the distribution. Indeed, the application of the scale parameter is pertinent when one wants to use the value of the mean time between events in the probability density function, as opposed to the total time until a given number of events occurs. 

That being said, in this document, as the two parameterizations only exist for the sake of convenience and are identical
in their results, only the parameterization 1 presented in \autoref{eq:gamma:parameterization1} will be considered and
used for proofs.

\subsection{Properties to prove}
\noindent Support: $x\in (0,\,\infty)$

\noindent Probability density function:
\begin{equation}\label{eq:gamma-pdf-2}
	f(x) = \frac{\lambda^\alpha}{\Gamma(\alpha)}x^{\alpha-1}e^{-\lambda x}
\end{equation}

\noindent Cumulative distribution function:
\begin{equation}\label{eq:gamma-cdf-2}
	F(x) = \frac{1}{\Gamma(\alpha)}\gamma(\alpha,\, \lambda x)
\end{equation}

\noindent Expected value, also known as the theoretical mean:
\begin{equation}\label{eq:gamma-expected-value-2}
	\mu = E(x) = \frac{\alpha}{\lambda}
\end{equation}

\noindent There is no simple closed form equation for the median of a gamma distribution.

\noindent Mode:
\begin{equation}\label{eq:gamma-mode-2}
	\text{Mode} = \frac{(\alpha-1)}{\lambda}\text{ for }\alpha \geq 1
\end{equation}

\noindent Variance:
\begin{equation}\label{eq:gamma-variance-2}
	\text{Var}(x) = \frac{\alpha}{\lambda^2}
\end{equation}

\noindent Skewness:
\begin{equation}\label{eq:gamma-skewness-2}
	\text{Skewness} = \frac{2}{\sqrt{\alpha}}
\end{equation}

\noindent Excess kurtosis:
\begin{equation}\label{eq:gamma-kurtosis-2}
	\text{Kurtosis} = \frac{6}{\alpha}
\end{equation}

\noindent Entropy:
\begin{equation}\label{eq:gamma-entropy-2}
	\text{Entropy} = \alpha + \ln{\lambda} + \ln{\Gamma(\alpha)} + (1-\alpha)\psi(\alpha)
\end{equation}

\noindent Moment generating function:
\begin{equation}\label{eq:gamma-mgf-2}
	M(t) = \left(1-\frac{t}{\lambda}\right)^{-\alpha}\text{ for }t < \lambda
\end{equation}

\noindent Characteristic function:
\begin{equation}\label{eq:gamma-cf-2}
	\text{CF} = \left(1-\frac{it}{\lambda}\right)^{-\alpha}
\end{equation}

\noindent Methods of moments:
\begin{equation}\label{eq:gamma-mom-2}
	\begin{split}
		\alpha	&=	\frac{E(X)^2}{\text{Var}(X)}\\
		\lambda	&=	\frac{E(X)}{\text{Var}(X)}
	\end{split}
\end{equation}

% END OF DOCUMENT =====================================================================================================
\pagebreak

\begin{appendix}
	\phantomsection\listoffigures
	\phantomsection\listoftables
\end{appendix}

\pagebreak\phantomsection\printbibliography[heading=bibintoc,title={References}]
\end{document}
